
\section{Appendix A. Mathematical Derivations}

\subsection{A.1 [CLS] Retention Proof}

\textbf{Claim.} Under score clamping $s_0 = 10^4$, token 0 is always in 
$\mathcal{T}$.

\textbf{Proof.} All scores except $s_0$ are in $(-\infty, 10^4)$. When we 
compute top-$k$ with $k \geq 1$, position 0 has the strictly largest 
score (assuming ties are broken in favor of earlier indices). Therefore 
$0 \in \mathcal{T}$. $\square$

\subsection{A.2 Causality Preserved}

\textbf{Claim.} For decoders, PTP attention satisfies causality.

\textbf{Proof.} For query $i$ and pruned key $j \in \mathcal{T}$:
\begin{itemize}
\item If original position $\mathcal{T}[b, j] > i$, then the causal mask is
  $-\infty$, and $P_{ij} = 0$.
\item Otherwise, $\mathcal{T}[b, j] \leq i$, and attention is allowed.

\end{itemize}
Therefore $O_i$ depends only on tokens with original position $\leq i$. 
$\square$

\subsection{A.3 Approximation Error}

\textbf{Claim.} The difference between PTP and dense attention output is:

\[
\|O - O_\text{dense}\|_F \leq 2 \max_j \|V_j\|_F \cdot \sum_{j \notin \mathcal{T}} P_j
\tag{A.1}
\]

\textbf{Proof sketch.} Write $O = \sum_{j \in \mathcal{T}} P_j V_j$ and 
$O_\text{dense} = \sum_{j=1}^L P_j V_j$. The difference is 
$\sum_{j \notin \mathcal{T}} P_j V_j$. Taking norms:

\[
\|O - O_\text{dense}\|_F = \left\|\sum_{j \notin \mathcal{T}} P_j V_j\right\|_F \leq \sum_{j \notin \mathcal{T}} P_j \|V_j\|_F
\]

Using $\|V_j\|_F \leq \max_j \|V_j\|_F$:

\[
\|O - O_\text{dense}\|_F \leq \max_j \|V_j\|_F \cdot \sum_{j \notin \mathcal{T}} P_j \qquad \square
\]

\subsection{A.4 Complexity}

\textbf{PTP attention:}

\begin{table}[H]
\centering
\small
\begin{tabular}{ll}
\toprule
\textbf{Operation} & \textbf{Complexity} \\
\midrule
Scoring & $\mathcal{O}(B L d \cdot 64)$ \\
Top-$k$ & $\mathcal{O}(B L \log k)$ \\
Gather & $\mathcal{O}(B H k d_h)$ \\
Attention & $\mathcal{O}(B H L k d_h)$ \\
\bottomrule
\end{tabular}
\end{table}

\textbf{Reduction factor} from dense: $\rho = k / L$.

\subsection{A.5 Optimal Keep Ratio}

Assume retained attention mass $m(\rho)$ satisfies $m(1) = 1$ and 
$m(\rho)$ is concave. Then task accuracy is approximately 
$\text{Acc}(\rho) \approx \text{Acc}(1) \cdot m(\rho)$.

We seek:

\[
\rho^* = \arg\max_\rho \left[\text{Acc}(1) m(\rho) - \lambda \rho\right]
\tag{A.2}
\]

Setting derivative to zero:

\[
\text{Acc}(1) m'(\rho) = \lambda
\tag{A.3}
\]

For $m(\rho) = \sqrt{\rho}$ (typical), $\rho^* = (\text{Acc}(1) / (2\lambda))^2$.

\subsection{A.6 Score Network Gradient}

For loss $\mathcal{L}$ and score $s_i = \sigma(g_i)$ where 
$g_i = \text{MLP}(x_i)$:

\[
\frac{\partial \mathcal{L}}{\partial g_i} = \frac{\partial \mathcal{L}}{\partial s_i} \cdot \sigma(g_i)(1 - \sigma(g_i))
\tag{A.4}
\]

The gradient flows through the top-$k$ operation only via the selected 
tokens (a straight-through estimator is implicitly used for the 
non-selected ones).

\subsection{A.7 Comparison to Random Pruning}

\textbf{Claim.} Random pruning of $k$ tokens has expected accuracy:

\[
\mathbb{E}[\text{Acc}_\text{random}] \approx \text{Acc}(1) \cdot \frac{k}{L}
\tag{A.5}
\]

\textbf{Observation.} PTP significantly exceeds this baseline.

\subsection{A.8 Numerical Stability}

Score clamping bounds $s_i \in [-10^4, 10^4]$, avoiding overflow in 
sigmoid (which is bounded anyway) and preserving numerical stability 
in \texttt{topk}.

